\documentclass[11pt]{article}

\usepackage[margin=1.05in]{geometry}
\usepackage[T1]{fontenc}
\usepackage{lmodern}
\usepackage{amsmath,amssymb,amsthm,mathtools}
\usepackage{array,booktabs,longtable}
\usepackage{xcolor}
\usepackage{microtype}
\usepackage{hyperref}
\hypersetup{
  colorlinks=true,
  linkcolor=blue!55!black,
  citecolor=blue!55!black,
  urlcolor=blue!55!black,
  pdftitle={Exclusion of the fixed-linear primitive cubic-pencil quartic row}
}

\newtheorem{theorem}{Theorem}
\newtheorem{proposition}[theorem]{Proposition}
\newtheorem{lemma}[theorem]{Lemma}
\newtheorem{corollary}[theorem]{Corollary}
\theoremstyle{remark}
\newtheorem{remark}[theorem]{Remark}

\newcommand{\A}{\mathbb A}
\newcommand{\C}{\mathbb C}
\newcommand{\PP}{\mathbb P}
\newcommand{\Jac}{\operatorname{Jac}}
\newcommand{\rank}{\operatorname{rank}}
\newcommand{\Sym}{\operatorname{Sym}}
\newcommand{\ord}{\operatorname{ord}}

\title{\bfseries Exclusion of the fixed-linear primitive
cubic-pencil quartic row\\[3pt]
\large A self-contained, AI-assisted research-program note}
\author{Alec Kriebel\\with substantial AI assistance}
\date{First repository release timestamp (UTC):
  \texttt{\detokenize{2026-07-25T23:14:47Z}}}

\begin{document}
\maketitle

\begin{center}
\fcolorbox{black!45}{yellow!8}{%
\begin{minipage}{0.91\textwidth}
\small
\textbf{Review, priority, and verification notice.}
This manuscript is not peer reviewed.  Exact computer algebra checks the
identities and finite case ledgers encoded in the accompanying files; it is
evidence, not peer review.  Substantial AI assistance was used in the
mathematical exploration, proof organization, software, hostile internal
audits, and exposition.  No claim of worldwide priority is made.  The result
excludes one frozen quartic leading-form row; it does not prove the full
degree-four case of the Jacobian Conjecture.
\end{minipage}}
\end{center}

\begin{abstract}
Let $F:\A^3_\C\to\A^3_\C$ be a degree-four Keller map.  We prove that
$F$ is an automorphism when the canonical invariant tuple of its quartic
homogeneous part is
\[
 (\rank JH_4,e,a,b,\delta,\nu)=(2,1,3,1,1,1).
\]
Equivalently, the leading part is an invariant fixed linear divisor times a
primitive cubic pencil.  A coefficient-independent normalization gives
$H_4=h(p,q,0)$ and an intrinsic horizontal/unique-vertical split.  A
homogeneous first-integral valuation settles the horizontal branch and
reduces the vertical branch to two cubic companions over a triple vertical
member.  Weighted Jacobian identities then eliminate all root, collision,
and lower-jet branches.  On the final branch, where a binary quadratic
$W_0$ is nonzero, the degree-six identity forces the cubic pencil onto the
exact nonminimal boundary.  Arbitrary quadratic and cubic lower jets are
retained throughout.
\end{abstract}

\section{Invariant statement and frozen context}

After translating source and target, write a degree-four polynomial map as
\begin{equation}
 F=LX+H_2+H_3+H_4, \qquad L\in\operatorname{GL}_3(\C),
 \label{eq:F}
\end{equation}
where $H_i$ is homogeneous of degree $i$ and $H_4\ne0$.  The map is
\emph{Keller} if $\det JF\in\C^\times$.

For a rank-two $H_4$, let
\[
 h=\gcd(H_{4,1},H_{4,2},H_{4,3}),\qquad e=\deg h,\qquad G=H_4/h.
\]
Put
\[
 K_G=\C(G_i/G_j:G_j\ne0)\subset M=\C(\PP^2),
\]
and let $E_G$ be the relative algebraic closure of $K_G$ in $M$.  The curve
with function field $E_G$ is rational, so $E_G=\C(p/q)$ for coprime
homogeneous forms of a common least degree $a$.  If the reduced projective
image of $G$ has degree $\delta$ and $\nu=[E_G:K_G]$, the induced map from
the pencil line is represented by a basepoint-free binary triple $A$ of
degree
\begin{equation}
 b=\delta\nu,\qquad e+ab=4.                           \label{eq:arithmetic}
\end{equation}
This defines the invariant tuple in the theorem.

\begin{theorem}[Fixed-linear primitive cubic-pencil exclusion]
\label{thm:main}
Let $F$ be as in \eqref{eq:F}, of exact total degree four, and suppose
\[
 (\rank JH_4,e,a,b,\delta,\nu)=(2,1,3,1,1,1).
\]
If $F$ is Keller, then $F$ is a polynomial automorphism.
\end{theorem}

\begin{corollary}
No degree-four Keller counterexample belongs to frozen row
\texttt{Q2-E1-A3-B1-D1-N1}.
\end{corollary}

The project's frozen denominator contains fourteen disjoint inclusive
leading-form rows.  At this release the internal status is
\[
 3/14\ \text{certified excluded},\qquad
 5/14\ \text{provisional},\qquad
 6/14\ \text{open}.
\]
This is bookkeeping context, not a theorem about the other rows.  The
universal total-degree floor remains four.

\section{Canonical normalization}

\begin{proposition}\label{prop:normal}
After independent invertible source and target changes,
\begin{equation}
 H_4=(hp,hq,0)^T,                                    \label{eq:normal}
\end{equation}
where $h$ is linear, $p,q$ are coprime nonproportional cubics, and
\begin{equation}
 \C(p/q)\text{ is relatively algebraically closed in }\C(\PP^2).
 \label{eq:minimal}
\end{equation}
The target normalization has three intrinsic rank-two charts and does not
divide by a frozen coefficient pivot.
\end{proposition}

\begin{proof}
The construction preceding \eqref{eq:arithmetic} gives the exact polynomial
factorization
\[
 H_4=hA(p,q).
\]
Indeed, the curve with function field $E_G$ is dominated by $\PP^2$.
Restriction to a general source line makes it dominated by $\PP^1$, hence
rational.  The primitive triples $G$ and $A(p,q)$ define the same
projective rational map.  The substituted triple is primitive because
$p,q$ are coprime and $A$ is basepoint-free.  Two primitive triples over
the UFD $\C[x,y,z]$ defining the same projective map differ by a unit,
which can be absorbed into $A$.

Here $b=1$, so
\[
 A(u,v)=\mathbf a\,u+\mathbf b\,v
\]
for linearly independent $\mathbf a,\mathbf b\in\C^3$.  Complete them to a
target basis and apply its inverse.  This proves \eqref{eq:normal}.
Explicitly, the three charts are given by the nonvanishing of
\[
 \Delta_{01}=a_0b_1-a_1b_0,\quad
 \Delta_{02}=a_0b_2-a_2b_0,\quad
 \Delta_{12}=a_1b_2-a_2b_1.
\]
Adjoining the unused standard basis vector gives determinant, up to sign,
equal to the selected minor.  These charts cover the rank-two condition and
refer to no first-nonzero coefficient of $H_4$.  Finally, $a=3$ is the
least pencil degree by the relative-closure construction, giving
\eqref{eq:minimal}.  Arbitrary lower jets are transported bijectively.
\end{proof}

\section{Weighted identities and first-integral descent}

Put
\[
 P=hp,\qquad Q=hq,\qquad D(R)=\Jac(P,Q,R),
\]
and introduce
\[
 \mathcal J(\tau)=L+\tau JH_2+\tau^2JH_3+\tau^3JH_4,\qquad
 E_j=[\tau^j]\det\mathcal J(\tau).
\]
The Keller condition gives $E_j=0$ for $j>0$.  If
$G_3=(H_3)_3$ and $G_2=(H_2)_3$, the zero third row of $JH_4$ gives
\begin{equation}
 E_8=D(G_3)=0.                                       \label{eq:E8}
\end{equation}
After $G_3=0$, the third rows of both $JH_4$ and $JH_3$ vanish, and
\begin{equation}
 E_7=D(G_2)=0.                                       \label{eq:E7zero}
\end{equation}

\begin{lemma}[Homogeneous first-integral descent]\label{lem:descent}
If $0\ne R\in\C[x,y,z]$ is homogeneous of degree $d$ and $D(R)=0$, then
\begin{equation}
 \frac{R^4}{P^d}=\mathcal R(q/p)
 \quad\text{for some }\mathcal R\in\C(t).             \label{eq:descent}
\end{equation}
\end{lemma}

\begin{proof}
The functions $P,Q$ are algebraically independent: $Q/P=q/p$ is
nonconstant on $\PP^2$, while $P$ supplies the independent scaling
parameter.  In characteristic zero,
$dP\wedge dQ\wedge dR=0$ makes $R$ algebraic over
$\C(P,Q)=\C(q/p,P)$.

The quotient $\Theta=R^4/P^d$ has degree zero and lies in
$M=\C(\PP^2)$.  With a scaling coordinate $s$ one has
$\C(x,y,z)=M(s)$, $P=s^4P_0$, and $R=s^dR_0$, so $P$ is
transcendental over $M$.  Clear denominators in an algebraic relation for
$\Theta$ over $\C(q/p)(P)$ and collect powers of $P$.  Transcendence makes
each coefficient vanish; hence $\Theta$ is algebraic over $\C(q/p)$.
Relative algebraic closure \eqref{eq:minimal} gives
$\Theta\in\C(q/p)$.
\end{proof}

\section{The horizontal/unique-vertical split}

Let $\mathcal P=\langle p,q\rangle$ and restrict to $h=0$:
\[
 \rho_h:\mathcal P\longrightarrow\C[x,y,z]_3/(h).
\]
Rank two means that no pencil member is divisible by $h$ and defines the
\emph{horizontal} branch.  Rank one gives a unique projective vertical
member.  Rank zero would make $h$ divide both generators, contrary to
coprimality.  This split is intrinsic under a change of pencil generators.

\subsection{The horizontal branch}

Horizontality says that $h$ divides neither $p$ nor $q-\lambda p$ for any
$\lambda\in\C$.  Factoring the numerator and denominator of $\mathcal R$
in \eqref{eq:descent} therefore gives
$v_h(\mathcal R(q/p))=0$.  Since $v_h(P)=1$,
\[
 4v_h(R)=d.
\]
There is no nonzero homogeneous first integral of degree two or three.
Equations \eqref{eq:E8} and \eqref{eq:E7zero} give
\[
 G_3=G_2=0.
\]
The third component of $F$ is linear.  After linear changes,
$F=(F_1,F_2,z)$; each plane fibre is Keller of degree at most four.
The bounded plane theorem used below makes each fibre an automorphism,
whence $F$ is injective and therefore an automorphism.

\subsection{The vertical branch}

Set $h=z$ and choose the unique vertical member as the first generator:
\begin{equation}
 p=z^m r,\qquad 1\le m\le3,\qquad z\nmid rq.          \label{eq:vertical}
\end{equation}
The integer $m=v_z(p)$ is intrinsic.

\section{The quadratic-component exit}

\begin{lemma}\label{lem:quadratic}
Let $F:\A^3_\C\to\A^3_\C$ be Keller of degree at most four.  If a
nonzero target-linear combination of its components has degree at most two,
then $F$ is an automorphism.
\end{lemma}

\begin{proof}
Call the combination $f$.  Its gradient never vanishes because $JF$ is
invertible.  Write $\nabla f=HX+b$, with $H$ constant symmetric.  If
$b\in\operatorname{im}H$, the equation $HX=-b$ gives a critical point.
Thus $b\notin(\ker H)^\perp$.  Choose $v\in\ker H$ with $b^Tv\ne0$.
In linear coordinates with $v$ as the third direction,
\[
 f=g(y_1,y_2)+\beta y_3,\qquad\beta\ne0.
\]
The triangular map $T=(y_1,y_2,f)$ and its inverse both have degree at
most two.  After the target change selecting $f$, put $G=F\circ T^{-1}$.
Then
\[
 G=(G_1,G_2,x_3),\qquad \deg G\le8.
\]
Each fibre $x_3=c$ is a plane etale polynomial map of maximum component
degree at most eight.

Vistoli states on journal p.~80 that an etale polynomial self-map of
$\A^2$ over the algebraically closed characteristic-zero field fixed on
p.~79 is an isomorphism when its degree is at most $12$
\cite[p.~80]{Vistoli}.  He attributes this theorem to Moh and notes that
Moh proves the result through degree $100$, whereas $12$ is all Vistoli
uses.  Hence every fibre of $G$ is an automorphism.  Equal images under
$G$ have equal third coordinates and then equal first two coordinates, so
$G$ is injective.  The injective-etale theorem quoted by Vistoli, or
Ax--Grothendieck together with etaleness, makes $G$ an automorphism.
\end{proof}

\begin{remark}
This uses Moh's bounded plane theorem in the exact degree-$\le12$ form
quoted by Vistoli.  It does not assume the unresolved plane Jacobian
Conjecture.
\end{remark}

\section{The vertical multiplicity lemma}

\begin{lemma}\label{lem:multiplicity}
For \eqref{eq:vertical}, the cubic kernel of \eqref{eq:E8} is
\[
\begin{array}{c|c}
 m&\ker\bigl(D:\C[x,y,z]_3\to\C[x,y,z]_8\bigr)\\ \hline
 1&0\\
 2&0\\
 3&\langle z^3,q\rangle .
\end{array}
\]
\end{lemma}

\begin{proof}
Apply Lemma~\ref{lem:descent} to a nonzero cubic $R$, and put
$\eta=\ord_\infty\mathcal R$, with a degree-$n$ polynomial having order
$-n$ at infinity.  If a prime $f$ occurs in $p$ with multiplicity $a$,
valuation of \eqref{eq:descent} gives
\begin{equation}
 4v_f(R)-3\bigl(a+\mathbf1_{f=z}\bigr)=a\eta.
 \label{eq:valuation}
\end{equation}
For $m=1$, the $z$-equation gives $\eta\equiv2\pmod4$.  Every prime of
the quadratic $r$ has multiplicity $a=1$ or $2$, but its equation requires
$a(3+\eta)\equiv0\pmod4$, impossible because
$3+\eta\equiv1\pmod4$.

For $m=2$, the $z$-equation is $4v_z(R)=9+2\eta$, an opposite-parity
contradiction.  For $m=3$ it is
\[
 4v_z(R)=12+3\eta.
\]
Since $0\le v_z(R)\le3$, the possibilities are
$(v_z(R),\eta)=(3,0)$ and $(0,-4)$.  The first gives
$R\sim z^3$.  In the second, $\mathcal R$ has no finite pole, since a
pole would give a negative valuation along a finite pencil fibre while the
left side of \eqref{eq:descent} is nonnegative there.  Thus $\mathcal R$ is
a degree-four polynomial.  A root of multiplicity $n<4$ would require
$4\mid na$ for every component multiplicity $a\in\{1,2,3\}$ of a cubic
fibre, impossible in total degree three.  Hence
$\mathcal R=c(t-\lambda)^4$ and unique factorization gives
$R\sim q-\lambda z^3$.
\end{proof}

The cases $m=1,2$, and $m=3$ with $G_3=0$ are settled by
Lemma~\ref{lem:quadratic}.  If $m=3$ and
$G_3=\alpha z^3+\beta q\ne0$, then $\beta=0$ gives the
\emph{vertical companion} $G_3=z^3$.  If $\beta\ne0$, replacing $q$ by
$G_3$ gives the \emph{nonvertical companion} $G_3=q$.  They cannot merge:
the unique vertical member is divisible by $z$, and $q$ is not.

\section{Companion gauges and the root atlas}

Write
\[
 H_3=(U,V,G_3)^T,\qquad H_2=(A,B,W)^T,\qquad
 \{f,g\}=f_xg_y-f_yg_x.
\]
For $G_3=q$, the complete weight-seven solve modulo legal target shears is
\begin{equation}
 U=dz^3,\qquad V=zW+fz^3.                             \label{eq:nonvertgauge}
\end{equation}
For $G_3=z^3$,
\[
 E_7=z^3\{q,4zW-3U\}=0,
\]
and the complete legal gauge is
\begin{equation}
 U=\frac43zW+\sigma q,\qquad [z^3]V=0.               \label{eq:vertgauge}
\end{equation}
Target shears remove the free $z^3$ coefficients of $q$ and $V$ while
only renaming unrestricted lower jets.  The split $\sigma=0$ or
$\sigma\ne0$ is exhaustive.

The nonzero binary cubic $q_0=q|_{z=0}$ has root partition
$1+1+1$, $2+1$, or $3$, represented by $xy(x-y)$, $x^2y$, and $x^3$.
On the triple-root locus, the full parabolic stabilizer gives the three
minimal charts
\begin{equation}
\begin{aligned}
 q_C&=x^3+y^2z+\alpha xz^2+\beta z^3,\\
 q_B&=x^3+xyz+\beta z^3,\\
 q_E&=x^3+yz^2.
\end{aligned}                                        \label{eq:tripleatlas}
\end{equation}
The fourth shape is $q\in\C[x,z]_3$ and is exactly the nonminimal
boundary.  For the vertical companion the term $\beta z^3$ is removable;
for the nonvertical companion it is retained.

\section{The fifteen terminal groups}

Let SF, D, and T denote squarefree, double-plus-simple, and triple root
types.  The following table is disjoint and exhaustive.

\begin{longtable}{@{}>{\bfseries\raggedright\arraybackslash}p{0.07\textwidth}
                         >{\raggedright\arraybackslash}p{0.55\textwidth}
                         >{\raggedright\arraybackslash}p{0.29\textwidth}@{}}
\toprule
ID & Predicate & Exit\\
\midrule
\endhead
H  & horizontal restriction rank two & homogeneous valuation, then plane
fibres\\
Z1 & vertical $m=1$ & $G_3=0$, Lemma~\ref{lem:quadratic}\\
Z2 & vertical $m=2$ & $G_3=0$, Lemma~\ref{lem:quadratic}\\
Z3 & $m=3$, $G_3=0$ & Lemma~\ref{lem:quadratic}\\
N1 & $G_3=q$, $q_0$ SF & nonvertical nontriple obstruction\\
N2 & $G_3=q$, $q_0$ D & nonvertical nontriple obstruction\\
N3 & $G_3=q$, $q_0$ T & three-chart nonvertical obstruction\\
V1 & $G_3=z^3$, $\sigma\ne0$, $q_0$ SF, $\ell=0$ &
zero-$\ell$ obstruction\\
V2 & same, $\ell\ne0$ & nonzero-$\ell$ obstruction\\
V3 & $G_3=z^3$, $\sigma\ne0$, $q_0$ D, $\ell=0$ &
zero-$\ell$ obstruction\\
V4 & same, $\ell\ne0$ & nonzero-$\ell$, including both collisions\\
V5 & $G_3=z^3$, $\sigma\ne0$, $q_0$ T, $\gamma\ne0$ &
direct $E_6$ contradiction\\
V6 & same, $\gamma=0$ & $E_6$ forces $\ell=0$, then $E_5,E_4$\\
V7 & $G_3=z^3$, $\sigma=0$, $W_0=0$ &
five-chart $E_6,E_5,E_4$ obstruction\\
V8 & $G_3=z^3$, $\sigma=0$, $W_0\ne0$ &
$E_6$ forces the nonminimal boundary\\
\bottomrule
\end{longtable}

Here $W_0=W|_{z=0}$.  If $\sigma\ne0$, restriction of $E_6$ gives
\begin{equation}
 -\sigma q_0\{q_0,W_0\}=0.                            \label{eq:Wrestrict}
\end{equation}
For SF or D the degree-two bracket kernel is zero, so
$W=z\ell+\omega z^2$.  For T, \eqref{eq:Wrestrict} gives
$W_0=\gamma x^2$.  Thus the table covers the entire vertical companion.

\subsection{Nonvertical terminals N1--N3}

With \eqref{eq:nonvertgauge}, for a noncube $q_0$, restrictions of
$E_6,E_5,E_4$ give
\[
 A_0=0,\qquad\overline L_1=0,\qquad
 A=\alpha z^2\ \text{or}\ B_0=0.
\]
The two full systems have constant pivot minors $-2^{19}$ and $-2^{11}$,
uniformly for $q_0=xy(x-y)$ and $q_0=x^2y$, and force
\[
 B=z(\ell_{31}x+\ell_{32}y+\beta z),\qquad
 \overline L_2=0,\qquad
 \overline L_3=(\ell_{31},\ell_{32}).
\]
The first two rows of $L$ are multiples of $z$, so $\det L=0$.

On the three charts \eqref{eq:tripleatlas}, the full $E_6,E_5$ systems have
constant $14\times14$ minors
\[
 -2^{24}3^8,\qquad -2^{18}3^6,\qquad -2^{20}3^7,
\]
and give the same dependent-row conclusion.  The moduli $\alpha,\beta$ are
not specialized.

\subsection{Vertical terminals V1--V6}

For SF or D and $\ell=ux+vy\ne0$, the complete binary $E_5$ kernel is
$V_0=\kappa q_0$, $\overline L_3=0$, except on the two D collision
lines, whose larger kernels are retained.  Raw coefficients give
\[
\begin{array}{c|cc}
\text{SF}&[x^4yz]E_6=\sigma u&[xy^4z]E_6=-\sigma v\\
\text{D, noncollision}&[x^4yz]E_6=\sigma u&
[x^3y^2z]E_6=-2\sigma v.
\end{array}
\]
On the collision kernels $\ell=cx$ and $\ell=cy$, the decisive
coefficients are $\sigma c$ and $-2\sigma c$.  Hence $\ell\ne0$ is
impossible.

If $\ell=0$, the complete degree-six solution is
\[
 V=kq+\frac z\sigma(A-a_5z^2)
   -\frac4{3\sigma}z^2(\ell_{31}x+\ell_{32}y).
\]
Constant pivots $2^5 3^{11}\sigma^8$ and
$-2^4 3^5\sigma^5$ prove completeness of the $E_6,E_5$ solves.
They give $\ell_{31}=\ell_{32}=0$, and $E_4$ gives
\[
 \ell_{21}=\frac{k}{\sigma}\ell_{11},\qquad
 \ell_{22}=\frac{k}{\sigma}\ell_{12}.
\]
Thus $\det L=0$.

For T, put $W=\gamma x^2+z(ux+vy+\omega z)$.  If $\gamma\ne0$, the
three charts give
\[
\begin{aligned}
[x^4yz]E_6&=4\gamma\sigma,\\
[x^5z]E_6&=\sigma(2\gamma-3v),&
-\tfrac16[x^3yz^2]E_6+[xy^2z^3]E_6
  &=-\tfrac{\sigma}{3}(\gamma+v),\\
\tfrac23[x^4z^2]E_6+[xyz^4]E_6&=\tfrac{10}{3}\gamma\sigma .
\end{aligned}
\]
Each chart is impossible.  If $\gamma=0$, every chart first gives
$[x^5z]E_6=-3\sigma v$, and a second chartwise linear combination gives a
nonzero constant times $\sigma u$.  Hence $u=v=0$, after which the complete
$E_6,E_5,E_4$ solution above makes $\det L=0$.

\subsection{The zero-parameter, zero-$W_0$ terminal V7}

Write $W=z(ux+vy+\omega z)$.  The SF, D, and three minimal T charts form
a complete atlas.  A constant $5\times5$ subsystem of $E_6$ gives on all
five charts
\[
\begin{aligned}
A_{20}&=\frac29u^2,&A_{11}&=\frac49uv,&A_{02}&=\frac29v^2,\\
\ell_{31}&=\frac{9A_{10}-4u\omega}{12},&
\ell_{32}&=\frac{9A_{01}-4v\omega}{12}.
\end{aligned}
\]
The nontriple charts and $q_C$ force $u=v=0$ through cubic monomials of
$E_5$.  The $q_B,q_E$ charts either do the same or leave one explicit
$u\ne0$ compatibility family.  On $u=v=0$, $E_5,E_4$ force
\[
 \ell_{11}=\ell_{12}=\ell_{31}=\ell_{32}=0,
\]
so $\det L=0$.  On either exceptional $u\ne0$ family, two or three raw
$E_4$ coefficients require incompatible values of a coefficient of $V$.

\section{The final $W_0\ne0$ elimination}

Assume $\sigma=0$, so $U=\frac43zW$.  Row multilinearity gives
\begin{equation}
 3E_6=z\Phi,\qquad
 \Phi=4W\{q,W\}+9z^2\{A,q\}+12z^3\{q,L_3\}.
 \label{eq:Phi}
\end{equation}
The variables $B,V,L_1,L_2$ cancel or do not occur.  The lowest
$z$-coefficient gives
\[
 \{q_0,W_0\}=0.
\]
For binary forms of degrees three and two, dehomogenization gives
$2F'G-3FG'=0$.  Thus $F^2/G^3$ is constant, and unique factorization
gives
\begin{equation}
 q_0=\kappa L^3,\qquad W_0=\gamma L^2,\qquad
 \kappa\gamma\ne0.                                   \label{eq:powers}
\end{equation}
After a binary change take $L=x$ and retain every coefficient:
\[
\begin{aligned}
q={}&\kappa x^3+z(\alpha x^2+\beta xy+\chi y^2)
 +z^2(\delta x+\epsilon y)+\phi z^3,\\
W={}&\gamma x^2+z(ux+vy)+\omega z^2,\\
A={}&a_{20}x^2+a_{11}xy+a_{02}y^2
 +z(a_{10}x+a_{01}y)+a_{00}z^2.
\end{aligned}
\]
Six raw coefficients of $\Phi$ suffice.  The first four give
\[
\chi=0,\quad
r=2\beta\gamma-3\kappa v=0,\quad
f=27\kappa a_{02}+2\beta\gamma v-6\kappa v^2=0,\quad
h=9a_{02}\beta-2\beta v^2=0.
\]
Put $g=9a_{02}-v^2$.  The division-free combinations
\[
 f-vr=3\kappa g,\qquad h-\beta g=-\beta v^2
\]
give $g=0$ and $\beta v^2=0$.  Multiplying $r=0$ by $v^2$ gives
$-3\kappa v^3=0$, hence
\[
 v=\beta=a_{02}=0.
\]
The remaining two coefficients reduce to
\[
 27\kappa a_{11}+8\gamma^2\epsilon=0,\qquad
 9a_{11}\epsilon=0.
\]
Multiplying the first by $\epsilon$ and subtracting three $\kappa$ times
the second gives $8\gamma^2\epsilon^2=0$.  Therefore
$\epsilon=a_{11}=0$, and
\[
 q=\kappa x^3+\alpha x^2z+\delta xz^2+\phi z^3
   \in\Sym^3\langle x,z\rangle.
\]
For a cubic pencil containing $z^3$, this is exactly the nonminimal
boundary $q\in\Sym^3\langle z,L\rangle$.  Its canonical pencil degree is
$a=1$, not $a=3$, so it reclassifies into a separate frozen row.  No point
of the minimal row survives.  This closes V8 and proves
Theorem~\ref{thm:main}.

\section{Frozen pivots, retained calculations, and verification boundary}

The frozen coefficient partition orders the fifteen quartic monomials
component by component and takes the first nonzero coefficient among
$c_0,\ldots,c_{44}$.  Proposition~\ref{prop:normal} and the proof above
are global on the invariant row and divide by none of these coefficients.
Every nonempty pivot \texttt{C00}--\texttt{C29} therefore enters the
theorem directly.  If the first pivot were \texttt{C30} or later, the
first two target components of $H_4$ would vanish, forcing
$\rank JH_4\le1$.  Hence \texttt{C30}--\texttt{C44} are empty.

The post-freeze bridge checker expands the proof into 48 disjoint route
atoms mapping uniquely to the fifteen terminal groups.  It checks all 45
pivots and includes leading-form witnesses for \texttt{C00}--\texttt{C29};
those witnesses are not asserted to be Keller completions.

Detailed derivations and separate hostile reconstructions are retained in
the repository files
\path{WORKING_HORIZONTAL_FIXED_LINEAR_CUBIC_PENCIL.md},
\path{vertical_locus/WORKING_VERTICAL_FIXED_LINEAR_CUBIC_PENCIL.md},
the branch lemmas under \path{vertical_locus/},
\path{vertical_locus/a0_w0_nonzero_attack/NOTE.md}, and
\path{../WORKING_QUADRATIC_COMPONENT_EXIT.md}.  The final bridge report is
\path{../taxonomy_freeze/audit_bridge_q2_e1_a3_b1_d1_n1_v1/REPORT.md}.
The aggregate entry point is \path{verify_all_strict.sh}.

These programs check encoded algebra, rank certificates, collision strata,
coverage, and deliberate fault injections.  Their ``hostile audit'' label
denotes an internal adversarial workflow, not independent human peer review.
The literature theorems used in Lemma~\ref{lem:quadratic} remain black
boxes.

\section{Prior art, attribution, and disclosure}

The external bounded-degree input is Moh's plane theorem in the
degree-$\le12$ form quoted by Vistoli.  Moh's cited paper is the 1983 paper
on configurations of roots \cite{Moh}.  We do not attribute the present
three-dimensional row decomposition to Moh, and we do not confuse the plane
theorem with the separate unpublished Moh--Sathaye computer calculation for
three-dimensional degree three discussed by Vistoli.

No claim is made that this row theorem is absent from all prior literature.
The project's source-specific search and internal provenance record are not
a worldwide priority determination.  A fuller account of the AI roles and
limitations is in \path{AI_DISCLOSURE.md}.  No external person was contacted
by an AI system on behalf of this project.

\begin{thebibliography}{9}

\bibitem{Moh}
T.-T. Moh,
\newblock On the Jacobian conjecture and the configurations of roots,
\newblock \emph{J. Reine Angew. Math.} \textbf{340} (1983), 140--212.
\newblock \url{https://doi.org/10.1515/crll.1983.340.140}.

\bibitem{Vistoli}
Angelo Vistoli,
\newblock The Jacobian conjecture in dimension 3 and degree 3,
\newblock \emph{J. Pure Appl. Algebra} \textbf{142} (1999), 79--89.
\newblock \url{https://doi.org/10.1016/S0022-4049(98)00040-1}.

\bibitem{BCW}
Hyman Bass, Edwin H. Connell, and David Wright,
\newblock The Jacobian conjecture: reduction of degree and formal expansion
of the inverse,
\newblock \emph{Bull. Amer. Math. Soc. (N.S.)} \textbf{7} (1982), 287--330.
\newblock
\url{https://www.ams.org/bull/1982-07-02/S0273-0979-1982-15032-7/}.

\bibitem{Stacks}
The Stacks Project,
\newblock Theorem 41.14.1, Tag 025G,
\newblock \url{https://stacks.math.columbia.edu/tag/025G}.

\end{thebibliography}

\end{document}
